% math.tex — basic gradient validation (Q1 generic, Q2 non-generic)
%
% Smoothness analysis of the EHZ capacity as a function of dual vertices.
% Refines [prop:capacity-piecewise-smooth] from library/src/algorithms/math.tex
% with per-orbit feasibility analysis and precise smoothness classification.
%
% Sources:
%   - Generated data: experiments/numerics/gradient/numerics/*.jsonl
%   - Developer notes: experiments/numerics/gradient/numerics/logbook.md
%   - KKT formulation: library/src/kkt/math.tex
%   - Existing proposition: library/src/algorithms/math.tex
%
% Structure:
%   §1 Setup and notation
%   §2 Per-orbit feasibility and smoothness (lem:orbit-feasibility-open, lem:per-orbit-smooth)
%   §3 Orbit contraction at feasibility boundary (lem:orbit-contraction)
%   §4 KKT sensitivity (lem:kkt-sensitivity)
%   §5 Capacity smoothness classification (prop:capacity-smoothness-classification)
%
% Labels defined here:
%   lem:orbit-feasibility-open, lem:per-orbit-smooth, lem:orbit-contraction,
%   lem:kkt-sensitivity, prop:capacity-smoothness-classification


% ── §1 Setup and notation ──

\subsection*{Setup and notation}

Let $K = \{x \in \mathbb{R}^4 : a_i^\top x \leq 1,\; i = 1, \ldots, F\}$
be a convex polytope with dual vertices
$a = (a_1, \ldots, a_F) \in (\mathbb{R}^4)^F$.
We treat~$a$ as the parameter and study how the EHZ capacity
$c(a)$ depends on it.

For a cyclic permutation $\sigma = (\sigma_1, \ldots, \sigma_m)$
of a subset of~$\{1, \ldots, F\}$,
the KKT system (see \texttt{library/src/kkt/math.tex}) is
\begin{equation}\label{eq:kkt-system}
  M_\sigma(a)\, x = b,
  \qquad
  M_\sigma = \begin{pmatrix}
    H & A & \mathbf{1} \\
    A^\top & 0 & 0 \\
    \mathbf{1}^\top & 0 & 0
  \end{pmatrix},
  \qquad
  x = \begin{pmatrix} \beta \\ \mu \\ \xi \end{pmatrix},
  \qquad
  b = \begin{pmatrix} 0 \\ 0 \\ 1 \end{pmatrix},
\end{equation}
where $H_{il} = \omega_0(a_{\sigma(l)}, a_{\sigma(i)})$ is the
symplectic form evaluated on the dual vertices indexed by~$\sigma$,
$A_{id} = (a_{\sigma(i)})_d$ is the coordinate matrix,
and $\mathbf{1} = (1, \ldots, 1)^\top$.
The matrix~$M_\sigma(a)$ depends smoothly on~$a$
($H$ is bilinear and $A$ is linear in the dual vertices).

When $M_\sigma(a)$ is non-singular, the KKT system has the
unique solution
\[
  x(a) = \bigl(\beta(\sigma; a),\; \mu(\sigma; a),\; \xi(\sigma; a)\bigr)
  = M_\sigma(a)^{-1}\, b.
\]
We write $\beta_k(\sigma; a)$ for the $k$-th component
($k = 1, \ldots, m$), and define
\[
  Q(\sigma; a) = \tfrac{1}{2}\,\beta(\sigma; a)^\top
  H(a)\,\beta(\sigma; a),
  \qquad
  A_\sigma(a) = \frac{1}{2\,Q(\sigma; a)}.
\]
An orbit~$\sigma$ is \emph{feasible} at~$a$ when
$M_\sigma(a)$ is non-singular,
$\beta(\sigma; a) \geq 0$ componentwise,
and $Q(\sigma; a) > 0$.
The EHZ capacity is
\begin{equation}\label{eq:capacity-min}
  c(a) = \min\bigl\{A_\sigma(a) :
    \sigma \text{ feasible at } a\bigr\}.
\end{equation}


% ── §2 Per-orbit feasibility and smoothness ──

\begin{unverified}
\begin{lemma}[Feasibility region is open]\label{lem:orbit-feasibility-open}
Fix a cyclic permutation~$\sigma$.
Suppose at~$a_0$ the KKT matrix~$M_\sigma(a_0)$ is non-singular
and the solution has $\beta(\sigma; a_0) > 0$ componentwise.
Then there exists $\varepsilon > 0$ such that for all
$\|a - a_0\| < \varepsilon$,
the matrix~$M_\sigma(a)$ is non-singular and
$\beta(\sigma; a) > 0$.
\end{lemma}

\begin{proof}
The map $a \mapsto \det M_\sigma(a)$ is continuous and non-zero at~$a_0$,
so $M_\sigma(a)$ is non-singular in a neighborhood of~$a_0$.
On this neighborhood, $a \mapsto M_\sigma(a)^{-1} b$ is smooth
(Cramer's rule: each entry of~$M^{-1}$ is a ratio of
polynomials in the entries of~$M$, with denominator $\det M \neq 0$).
Since~$\beta(\sigma; a_0) > 0$ componentwise and~$\beta(\sigma; a)$ is
continuous, $\beta(\sigma; a) > 0$ for $a$ sufficiently close to~$a_0$.
\end{proof}
\end{unverified}


\begin{unverified}
\begin{lemma}[Per-orbit action is smooth]\label{lem:per-orbit-smooth}
Under the hypotheses of Lemma~\ref{lem:orbit-feasibility-open},
and assuming additionally that $Q(\sigma; a_0) > 0$,
the per-orbit action $A_\sigma(a) = 1/(2\,Q(\sigma; a))$
is $C^\infty$ in~$a$ on the neighborhood where~$M_\sigma(a)$ is
non-singular and $\beta(\sigma; a) > 0$.
\end{lemma}

\begin{proof}
$\beta(\sigma; a)$ is smooth in~$a$ (smooth entries composed with matrix inversion).
$H(a)$ is smooth in~$a$ (bilinear in dual vertices).
Hence $Q(\sigma; a) = \tfrac{1}{2}\beta^\top H\,\beta$
is smooth as a composition of smooth maps.
Since $Q(\sigma; a_0) > 0$ and~$Q$ is continuous,
$Q(\sigma; a) > 0$ on a neighborhood, and
$A_\sigma(a) = 1/(2Q)$ is smooth there.
\end{proof}
\end{unverified}



% ── §3 Orbit contraction at feasibility boundary ──

\begin{unverified}
\begin{lemma}[Orbit contraction at feasibility boundary]%
\label{lem:orbit-contraction}
Let $\sigma = (\sigma_1, \ldots, \sigma_m)$ be a cyclic permutation
and suppose the KKT system for~$\sigma$ at~$a$ has a solution
$(\beta, \mu, \xi)$ with $\beta_k = 0$ for some index~$k$
and $\beta_j > 0$ for all $j \neq k$.
Let $\sigma' = (\sigma_1, \ldots, \sigma_{k-1}, \sigma_{k+1}, \ldots, \sigma_m)$
be the contracted permutation.
Then:
\begin{enumerate}
\item[\textup{(a)}]
  $(\beta', \mu, \xi)$ where $\beta' = (\beta_1, \ldots, \beta_{k-1},
  \beta_{k+1}, \ldots, \beta_m)$ satisfies the KKT system for~$\sigma'$ at~$a$.
\item[\textup{(b)}]
  $A_\sigma(a) = A_{\sigma'}(a)$.
\end{enumerate}
\end{lemma}

\begin{proof}
Write the KKT system for~$\sigma$ in block rows:
\[
  \text{(I)}\ (H\beta)_j + (A\mu)_j + \xi = 0 \;\; \forall j, \qquad
  \text{(II)}\ A^\top \beta = 0, \qquad
  \text{(III)}\ \mathbf{1}^\top \beta = 1.
\]
Let $S' = \{1,\ldots,m\} \setminus \{k\}$ index the positions
of~$\sigma'$.
Let $H'$ denote the $(m{-}1) \times (m{-}1)$ submatrix of~$H$
obtained by deleting row and column~$k$,
and $A'$ the $(m{-}1) \times 4$ submatrix of~$A$
obtained by deleting row~$k$.

\emph{(a).}
For $j \in S'$:
$(H\beta)_j = \sum_{l \in S'} H_{jl} \beta_l + H_{jk} \cdot 0
= (H'\beta')_j$.
Similarly $(A\mu)_j = (A'\mu)_j$.
So equation~(I) for $j \in S'$ gives (I) for the $\sigma'$-system.
For (II): $A^\top \beta = \sum_{j \in S'} \beta_j a_{\sigma(j)} = {A'}^\top \beta' = 0$.
For (III): $\mathbf{1}^\top \beta' = \mathbf{1}^\top \beta - 0 = 1$.

\emph{(b).}
By Lemma~\ref{lem:well-defined},
$\beta^\top H\,\beta = -\xi$ for any solution of the
$\sigma$-system. Since $(\beta', \mu, \xi)$ satisfies
the $\sigma'$-system by~(a), the same lemma gives
${\beta'}^\top H'\,\beta' = -\xi$.
Hence $Q_\sigma = Q_{\sigma'} = -\xi/2$ and
$A_\sigma = 1/(2Q) = A_{\sigma'}$.
\end{proof}
\end{unverified}


% ── §4 KKT sensitivity ──

\begin{unverified}
\begin{lemma}[Sensitivity of the KKT solution]\label{lem:kkt-sensitivity}
Fix~$\sigma$ and suppose $M_\sigma(a_0)$ is non-singular
with solution $x_0 = M_\sigma(a_0)^{-1} b$.
Then $a \mapsto x(\sigma; a) = M_\sigma(a)^{-1} b$
is smooth near~$a_0$, and its derivative with respect to
the dual vertex~$a_j$ is
\begin{equation}\label{eq:kkt-sensitivity}
  \frac{\partial x}{\partial a_j}
  = -M_\sigma(a_0)^{-1}\,
    \frac{\partial M_\sigma}{\partial a_j}\, x_0.
\end{equation}
In particular, $\partial\beta_k(\sigma; a_0) / \partial a_j$
is the $k$-th component of this vector.
\end{lemma}

\begin{proof}
Since $M_\sigma(a_0)$ is non-singular and $\det M_\sigma(a)$ is continuous,
$M_\sigma(a)$ remains non-singular near~$a_0$,
so $x(a) = M_\sigma(a)^{-1} b$ is smooth.
Differentiating $M_\sigma(a)\,x(a) = b$ at~$a_0$
and using $\partial b / \partial a_j = 0$
(since $b = (0, 0, 1)$ is independent of~$a$):
\[
  \frac{\partial M_\sigma}{\partial a_j}\, x_0
  + M_\sigma(a_0)\, \frac{\partial x}{\partial a_j} = 0.
\]
Solving gives~\eqref{eq:kkt-sensitivity}.
\end{proof}
\end{unverified}


% ── §5 Capacity smoothness classification ──

\begin{unverified}
\begin{proposition}[Smoothness classification of $c(a)$]%
\label{prop:capacity-smoothness-classification}
\begin{enumerate}
\item[\textup{(a)}]
  \emph{Generic points.}
  If at~$a_0$ a unique orbit~$\sigma^*$ achieves the minimum
  in~\eqref{eq:capacity-min}, with $\beta(\sigma^*; a_0) > 0$,
  and every other feasible orbit has action strictly above~$c(a_0)$,
  then $c$ is $C^\infty$ in a neighborhood of~$a_0$,
  with gradient given by the envelope theorem formula
  (Lemma~\ref{lem:cap-derivative}).

\item[\textup{(b)}]
  \emph{Switching boundary (interior orbits).}
  If $r \geq 2$ orbits $\sigma_1, \ldots, \sigma_r$ achieve the
  minimum, each with $\beta(\sigma_i; a_0) > 0$,
  with distinct gradients
  ($\nabla_a A_{\sigma_i} \neq \nabla_a A_{\sigma_j}$
  for some~$i, j$),
  and every other feasible orbit has action strictly above~$c(a_0)$,
  then $c$ is Lipschitz but not differentiable at~$a_0$.
  The directional derivative is
  $D_d\, c(a_0) = \min_{i=1,\ldots,r} \nabla_a A_{\sigma_i}(a_0) \cdot d$.

% [GAP - degenerate tie case: r orbits tied with matching gradients.
%  c is C^1 there but possibly not C^2. Non-generic (higher codimension).
%  Not formalized.]

\item[\textup{(c)}]
  \emph{Genericity.}
  The switching locus (where $\geq 2$ orbits tie)
  has measure zero in the space of dual vertices.
\end{enumerate}
\end{proposition}

\begin{proof}
\emph{(a).}
By Lemma~\ref{lem:per-orbit-smooth},
$A_{\sigma^*}$ is $C^\infty$ near~$a_0$.
% [TODO: JÖRN - the gap argument assumes that every competing orbit's
%  action A_sigma(a) is continuous near a_0. This holds when each
%  competing sigma has non-singular M(a) (by lem:per-orbit-smooth),
%  but could fail for orbits with singular KKT matrix.
%  It also does not account for orbits that are infeasible at a_0
%  but become feasible nearby with lower action (orbit appearance).
%  Both gaps are flagged as open questions at the end of this file.]
Assuming each competing orbit's action is continuous near~$a_0$
and no new orbits appear:
the gap $\delta = \min_{\sigma \neq \sigma^*}
A_\sigma(a_0) - A_{\sigma^*}(a_0) > 0$ is preserved
in a neighborhood, so $c(a) = A_{\sigma^*}(a)$ there.

\emph{(b).}
Each~$A_{\sigma_i}$ is $C^\infty$ near~$a_0$
(by Lemma~\ref{lem:per-orbit-smooth},
since $\beta(\sigma_i; a_0) > 0$).
Write $g_i = \nabla_a A_{\sigma_i}(a_0)$.
Since each $A_{\sigma_i}$ is $C^2$, Taylor's theorem gives
$|A_{\sigma_i}(a_0 + td) - A_{\sigma_i}(a_0) - t\, g_i \cdot d|
\leq C_i\, t^2$ for small~$t$ and some constants~$C_i$.
Set $C = \max_i C_i$.

\emph{Upper bound.}
Each $\sigma_i$ remains feasible near~$a_0$
(by Lemma~\ref{lem:orbit-feasibility-open}), so
$c(a_0 + td) \leq A_{\sigma_i}(a_0 + td)$ for all~$i$.
Hence $c(a_0 + td) \leq \min_i A_{\sigma_i}(a_0 + td)
\leq c(a_0) + t\,\min_i(g_i \cdot d) + Ct^2$.

\emph{Lower bound.}
By hypothesis, every feasible orbit outside
$\{\sigma_1, \ldots, \sigma_r\}$ has action strictly
above~$c(a_0)$.
By continuity, these orbits' actions remain
above $c(a_0) + \delta/2$ for some $\delta > 0$
and small~$t$, so they cannot undercut the minimum.
Hence $c(a_0 + td) = \min_i A_{\sigma_i}(a_0 + td) + O(t^2)
\geq c(a_0) + t\,\min_i(g_i \cdot d) - Ct^2$.

Combining: $D_d\, c(a_0) = \min_i (g_i \cdot d)$.

\emph{Non-differentiability.}
Suppose for contradiction that $c$ is differentiable at~$a_0$
with gradient~$g$.
Then $g \cdot d = D_d c = \min_i (g_i \cdot d)$ for all~$d$.
A function that is both linear ($d \mapsto g \cdot d$)
and the pointwise minimum of linear functions
($d \mapsto \min_i(g_i \cdot d)$) requires all $g_i$ to
agree: taking $d = g - g_i$ gives $-\|g - g_i\|^2 \geq 0$.
This contradicts the distinct gradients hypothesis.

\emph{(c).}
$A_{\sigma_1}(a) = A_{\sigma_2}(a)$ is one equation in $4F$ parameters.
% [TODO: JÖRN - to conclude codimension 1, need
%  nabla_a(A_{sigma_1} - A_{sigma_2}) != 0 generically
%  (transversality). Plausible but not proved here.]
Assuming the two gradients are distinct on the tie locus,
the implicit function theorem gives a codimension-$1$
hypersurface. Finitely many orbit pairs, so the union has
measure zero.
\end{proof}
\end{unverified}
